% Trénování po batchích a epochách jako vývojový diagram
% (makra \flow... z assets/macros.tex).
\begin{tikzpicture}[
    x=1cm,
    y=1cm,
    font=\small
]

    \flowio{data}{(0,0)}{trénovací\\množina \(D\)}
    \flowstep{split}{(2.3,0)}{rozdělení\\na batche}
    \flowio{batch1}{(4.7,1.0)}{\(\mathcal{B}_1\)}
    \flowio{batch2}{(4.7,0)}{\(\mathcal{B}_2\)}
    \flowio{batchs}{(4.7,-1.7)}{\(\mathcal{B}_s\)}
    \flowstep{train}{(7.3,0)}{průchod sítí}
    \flowstep{update}{(10.0,0)}{úprava\\parametrů}

    \flowarrow{data}{split}
    \flowarrow{train}{update}
    \flowarrow{split.east}{batch1.west}
    \flowarrow{split.east}{batch2.west}
    \flowarrow{split.east}{batchs.west}

    \flowarrow{batch1.east}{train.west}
    \flowarrow{batch2.east}{train.west}
    \flowarrow{batchs.east}{train.west}

    \node at (4.7,-0.75) {\(\vdots\)};

    \draw[flowarrow]
        (update.north) -- ++(0,1.1) -- ++(-10,0) node[flowlabel, midway, above] {další epocha} -- (data.north);

\end{tikzpicture}
